\documentclass[11pt,a4paper]{article}

\usepackage[T1]{fontenc}
\usepackage[utf8]{inputenc}
\usepackage[spanish,es-noshorthands,es-tabla]{babel}
\usepackage{lmodern}
\usepackage{microtype}
\usepackage[margin=2.3cm]{geometry}
\usepackage[table]{xcolor}
\usepackage{booktabs}
\usepackage{tabularx}
\usepackage{array}
\usepackage{multirow}
\usepackage{enumitem}
\usepackage{tcolorbox}
\usepackage{titlesec}
\usepackage{amsmath,amssymb}
\usepackage{graphicx}
\usepackage{caption}
\usepackage{tikz}
\usepackage{pgfplots}
\pgfplotsset{compat=1.16}
\usepackage{tocloft}
\setlength{\cftsecnumwidth}{2em}
\renewcommand{\cftsecleader}{\cftdotfill{\cftdotsep}}
\usepackage[hidelinks]{hyperref}

\usetikzlibrary{babel,arrows.meta,positioning,calc,shapes.geometric}
\tcbuselibrary{skins,breakable}

\definecolor{azulnuevo}{HTML}{1F4E79}
\definecolor{grisrep}{HTML}{6B7280}
\definecolor{ambar}{HTML}{B45309}
\definecolor{verde}{HTML}{15803D}
\definecolor{fondonuevo}{HTML}{F3F7FB}
\definecolor{fondorep}{HTML}{F4F4F5}
\definecolor{fondoaviso}{HTML}{FFF7ED}
\definecolor{fondoej}{HTML}{F0FDF4}
\definecolor{critrojo}{HTML}{DC2626}
\definecolor{azultemp}{HTML}{1D4ED8}
\definecolor{rojotard}{HTML}{B91C1C}

\titleformat{\section}{\Large\bfseries\color{azulnuevo}}{\thesection}{0.6em}{}
\titleformat{\subsection}{\large\bfseries}{\thesubsection}{0.6em}{}
\titleformat{\subsubsection}{\normalsize\bfseries}{\thesubsubsection}{0.6em}{}
\titlespacing*{\section}{0pt}{12pt}{6pt}
\titlespacing*{\subsection}{0pt}{9pt}{4pt}

\newtcolorbox{defbox}[1]{
  breakable, enhanced, colback=fondonuevo, colframe=azulnuevo,
  boxrule=0.5pt, left=6pt, right=6pt, top=4pt, bottom=4pt,
  title={\bfseries #1}, coltitle=white, colbacktitle=azulnuevo, fonttitle=\small}
\newtcolorbox{yavisto}{
  breakable, enhanced, colback=fondorep, colframe=grisrep,
  boxrule=0.4pt, left=6pt, right=6pt, top=4pt, bottom=4pt, fontupper=\small}
\newtcolorbox{aviso}[1]{
  breakable, enhanced, colback=fondoaviso, colframe=ambar,
  boxrule=0.5pt, left=6pt, right=6pt, top=4pt, bottom=4pt,
  title={\bfseries #1}, coltitle=white, colbacktitle=ambar, fonttitle=\small}
\newtcolorbox{ejemplo}[1]{
  breakable, enhanced, colback=fondoej, colframe=verde,
  boxrule=0.5pt, left=6pt, right=6pt, top=4pt, bottom=4pt,
  title={\bfseries #1}, coltitle=white, colbacktitle=verde, fonttitle=\small}
\newtcolorbox{enresumen}{
  breakable, enhanced, colback=white, colframe=azulnuevo,
  boxrule=0.4pt, leftrule=3pt, left=6pt, right=6pt, top=4pt, bottom=4pt,
  title={\bfseries En resumen}, coltitle=azulnuevo, colbacktitle=white,
  fonttitle=\small, attach title to upper={\ }, fontupper=\small}

\setlist[itemize]{itemsep=1pt, topsep=2pt, parsep=0pt}
\setlist[enumerate]{itemsep=1pt, topsep=2pt, parsep=0pt}
\captionsetup{font=small, labelfont=bf}
\setlength{\parskip}{3pt}
\setlength{\parindent}{0pt}

\tikzset{
  caja/.style={draw=azulnuevo, fill=fondonuevo, rounded corners=2pt, align=center, font=\small, inner sep=4pt},
  flecha/.style={-{Stealth[length=2.2mm]}, thick, draw=azulnuevo},
  tarea/.style={draw=black!80, thick, minimum width=17mm, minimum height=8mm, fill=white, font=\small},
  crit/.style={tarea, draw=critrojo, fill=critrojo!15, very thick},
  temp/.style={text=azultemp, font=\scriptsize\bfseries},
  tard/.style={text=rojotard, font=\scriptsize\bfseries},
  dep/.style={-{Stealth[length=2.2mm]}, thick, draw=azulnuevo!75},
  depc/.style={-{Stealth[length=2.2mm]}, very thick, draw=critrojo}
}
\newcommand{\tnode}[9][tarea]{\node[#1,
  label={[temp,anchor=south west,inner sep=1pt]north west:#6},
  label={[temp,anchor=south east,inner sep=1pt]north east:#7},
  label={[tard,anchor=north west,inner sep=1pt]south west:#8},
  label={[tard,anchor=north east,inner sep=1pt]south east:#9}] (#2) at (#3) {\textbf{#4}\hspace{5mm}#5};}

\definecolor{fondoenun}{HTML}{F8FAFC}
\newtcolorbox{enunciado}[1]{
  breakable, enhanced, colback=fondoenun, colframe=black!55,
  boxrule=0.5pt, left=6pt, right=6pt, top=4pt, bottom=4pt,
  title={\bfseries #1}, coltitle=white, colbacktitle=black!65, fonttitle=\small}
\newtcolorbox{plantilla}[1]{
  breakable, enhanced, colback=fondorep, colframe=grisrep,
  boxrule=0.5pt, left=6pt, right=6pt, top=4pt, bottom=4pt,
  title={\bfseries Dado en el enunciado: #1}, coltitle=white, colbacktitle=grisrep, fonttitle=\small, fontupper=\small}
\newtcolorbox{resultado}[1]{
  breakable, enhanced, colback=fondoej, colframe=verde,
  boxrule=0.8pt, left=6pt, right=6pt, top=4pt, bottom=4pt,
  title={\bfseries #1}, coltitle=white, colbacktitle=verde, fonttitle=\small}
\newcommand{\hueco}[1][1.6cm]{\rule{0pt}{#1}}
\title{\vspace{-1.6cm}\textbf{Respuestas de la guía de ejercicios}\\[4pt]
\large Resoluciones paso a paso, para comparar después de resolver}
\author{\normalsize Administración de Proyectos de Software (7502) --- UNS}
\date{\normalsize Enunciados en \texttt{ejercicios-examen.pdf}}

\begin{document}
\maketitle
\thispagestyle{empty}

\begin{yavisto}
\textbf{Cómo usar este documento.} Resolvé primero cada ejercicio con la guía y después compará. Cada respuesta
indica qué teoría interviene (con la sección del resumen final), la estrategia, el desarrollo y el resultado en
una caja verde. Todos los cálculos numéricos (redes, Puntos de Función, COCOMO, valor ganado, matrices) se
verificaron con un script de Python. El script de redes reproduce exactamente la resolución de la cátedra del
ejercicio 3 del TP3, así que sigue la misma convención. En los ejercicios abiertos (requerimientos, riesgos,
asignaciones) la respuesta es \emph{una} solución válida: lo que se evalúa es que esté justificada.
\end{yavisto}

\tableofcontents
\newpage

% =====================================================================
\section{Bloque A --- Scrum}
% =====================================================================
\subsection*{Ejercicio A1}
\textbf{Teoría}: resumen final §2 (historias, INVEST, criterios de aceptación, velocity).

\paragraph{a) Descripción.} Es un sistema de \textbf{reservas de un club deportivo}: los socios reservan, cancelan
y pagan canchas (pádel), reciben recordatorios; los administrativos ven la ocupación; los profesores publican sus
clases. Los ítems técnicos (migrar la base, backups, corregir duplicados) indican que el sistema ya está en
producción.

\paragraph{b) Historias de usuario vs.\ ítems técnicos.} Pregunta guía: ¿lo podría haber pedido un usuario sin
conocimientos técnicos, porque necesita hacer algo nuevo?
\begin{itemize}
  \item \textbf{Historias de usuario}: \#1, \#2, \#4, \#6, \#8 y \#9. Siguen el formato ``como [rol], quiero\ldots''
  y su beneficiario es un usuario (socio, administrativo, profesor).
  \item \textbf{Ítems del equipo}: \#3 (migrar a PostgreSQL: infraestructura invisible para el usuario), \#5
  (corregir un error: el usuario ve el efecto, pero el ítem es una acción técnica de diagnóstico y corrección) y
  \#7 (backups: operación). Están en infinitivo y su beneficiario directo es el sistema.
\end{itemize}

\paragraph{c) Selección con velocity 21.}
\begin{center}\small
\begin{tabular}{@{}c l c c c c c@{}}
\toprule
Orden & \# & Valor & Esfuerzo & Ratio & Acumulado & ¿Entra?\\
\midrule
1 & 2 cancelar reserva & 5 & 2 & 2,50 & 2 & sí\\
2 & 1 reservar cancha & 8 & 5 & 1,60 & 7 & sí\\
3 & 5 corregir duplicados & 4 & 3 & 1,33 & 10 & sí\\
4 & 6 recordatorio & 3 & 3 & 1,00 & 13 & sí\\
5 & 4 ocupación semanal & 6 & 8 & 0,75 & 21 & sí\\
6 & 8 pagar con tarjeta & 7 & 13 & 0,54 & 34 & no\\
7 & 7 backups & 2 & 5 & 0,40 & 26 & no\\
7 & 9 horarios de clases & 2 & 5 & 0,40 & 26 & no\\
9 & 3 migrar a PostgreSQL & 3 & 8 & 0,38 & 29 & no\\
\bottomrule
\end{tabular}
\end{center}
\begin{resultado}{Resultado}
Entran \textbf{\#2, \#1, \#5, \#6 y \#4}, con \textbf{21 puntos}: la capacidad queda usada exactamente. Ningún otro
ítem cabe.
\end{resultado}
\textbf{Discusión}: \#3 (migración) y \#7 (backups) quedaron afuera por su ratio, pero un PO razonable podría
subir los backups (riesgo de perder datos) aunque su ratio sea bajo: el ratio orienta, no decide. Si se cambiara
\#4 (8 puntos) por \#7 (5 puntos), sobrarían 3 puntos.

\paragraph{d) Criterios de aceptación.}
\textbf{\#2 Cancelar reserva}: el socio puede ver sus reservas futuras; puede seleccionar una y cancelarla; si
faltan menos de X horas para el turno, el sistema informa que no se puede cancelar sin cargo (regla a acordar con
el PO); si la cancelación es válida, la cancha queda libre para otros socios; luego se muestra la lista de
reservas actualizada y se confirma la cancelación.
\textbf{\#6 Recordatorio}: el socio con una reserva recibe un recordatorio el día anterior, con cancha, día y
hora; si la reserva se canceló, no se envía; si el socio desactivó los recordatorios, no se envía; el
recordatorio incluye un enlace para cancelar. Patrón: lo que puede hacer el usuario, reglas ``si\ldots
entonces'', confirmación visible.

\paragraph{e) Agregar \#8 a mitad del sprint.} No se puede: una vez que empieza el sprint \textbf{el alcance
queda congelado}. El ítem va al Product Backlog y el PO lo prioriza para el próximo sprint (además, con 13
puntos ni siquiera entraría en la velocity). Si fuera realmente urgente, la única salida en Scrum es cancelar el
sprint, algo excepcional.

% =====================================================================
\section{Bloque B --- Requerimientos y riesgos}
% =====================================================================
\subsection*{Ejercicio B1: turnos en centros de salud municipales}
\textbf{Teoría}: resumen final §5 (requerimientos, atributos, soluciones disfrazadas) y §6 (riesgos).
\textbf{Estrategia}: subrayar el enunciado buscando cifras, plazos, actores, restricciones y faltantes (es lo que
hace la cátedra en la resolución del caso de salud), y después recorrer las categorías de riesgo de Pressman
preguntando cuáles aplican.

\paragraph{a) Requerimientos funcionales.}
\begin{enumerate}[label=F\arabic*]
  \item El paciente puede solicitar un turno por web o app, eligiendo centro, especialidad, profesional y horario disponible.
  \item El paciente puede consultar y cancelar sus turnos.
  \item El administrativo puede otorgar, consultar y cancelar turnos en mostrador.
  \item El director de cada centro define la agenda semanal de cada profesional.
  \item El sistema bloquea la agenda por feriados y licencias.
  \item El sistema valida que el paciente esté en el padrón provincial de beneficiarios antes de otorgar un turno.
  \item El sistema envía un recordatorio por SMS 48 horas antes del turno.
  \item El paciente puede anotarse en una lista de espera; cuando se cancela un turno, el sistema lo ofrece al
  primero de la lista.
  \item El sistema registra si el paciente asistió (turno usado o perdido).
  \item El sistema emite un reporte mensual de turnos dados, usados y perdidos por centro y especialidad, con el
  porcentaje de ausentismo.
\end{enumerate}
\begin{center}\small
\begin{tabularx}{\textwidth}{@{}l X X@{}}
\toprule
Campo & F1 & F8\\
\midrule
categoría & funcional & funcional\\
descripción & El paciente registrado solicita un turno eligiendo centro, especialidad, profesional, día y
horario entre los disponibles; el sistema lo reserva y lo muestra. & Si no hay turnos, el paciente se anota en
la lista de espera de una especialidad y centro; al cancelarse un turno, el sistema se lo ofrece al primero de
la lista por SMS, con un plazo para aceptarlo.\\
términos & paciente, turno, agenda, centro, especialidad & lista de espera, oferta de turno, plazo de aceptación\\
justificación & reducir la espera telefónica de más de 40 minutos & aprovechar los turnos cancelados (25\% de
ausentismo)\\
prioridad & alta & media\\
dependencias & F4 (agenda), F6 (padrón) & F1, F2, F7 (SMS)\\
factibilidad & existen soluciones maduras de agendas en línea & requiere un proveedor de SMS\\
verificación & un paciente empadronado reserva un turno libre y el horario deja de estar disponible para otros &
al cancelar un turno, el primero de la lista recibe la oferta y, si acepta, el turno queda a su nombre\\
\bottomrule
\end{tabularx}
\end{center}

\paragraph{b) No funcionales y atributo crítico.} Usabilidad para adultos mayores; tiempo de respuesta;
disponibilidad (la web y la app deben estar operativas, y los centros rurales deben poder seguir trabajando con
conectividad irregular); seguridad y confidencialidad de los datos de salud (cumplimiento de la ley de datos
personales); capacidad (85.000 consultas por año, picos en la campaña de vacunación); restricción de diseño:
web, app móvil y mostrador; fecha (8 meses) y presupuesto fijo.
\begin{center}\small
\begin{tabularx}{\textwidth}{@{}l X@{}}
\toprule
\multicolumn{2}{@{}l}{\textbf{Atributo: tiempo para sacar un turno (usabilidad)}}\\
\midrule
Escala & minutos desde que el paciente inicia la solicitud hasta que recibe la confirmación del turno\\
Test & prueba con 20 pacientes mayores de 65 años sin capacitación, en la app y en la web; se toma el promedio\\
Peor caso & 10 minutos (por encima, la gente seguiría llamando por teléfono y el sistema no resuelve el problema)\\
Nivel planificado & 3 minutos\\
Mejor caso & 1 minuto\\
Nivel actual & más de 40 minutos de espera telefónica en horas pico (dato de la Secretaría)\\
\bottomrule
\end{tabularx}
\end{center}
Otro atributo crítico válido: ausentismo (\% de turnos perdidos), con nivel actual 25\%, planificado 10\% y peor
caso 20\%; o disponibilidad del sistema en horario de atención.

\paragraph{c) Solución disfrazada y requerimiento ambiguo.}
\begin{itemize}
  \item \textbf{Solución disfrazada}: ``un sistema de turnos como el que usa la ciudad vecina''. Es una solución
  elegida de antemano. Se pregunta \emph{¿por qué ese?}: probablemente quiere reducir la espera y el ausentismo
  con algo ya probado. Ese es el requerimiento real; el sistema de la ciudad vecina es una solución candidata
  (\emph{non-sancta}) que se evalúa contra los atributos (¿consulta el padrón provincial?, ¿funciona con mala
  conectividad?, ¿cuesta lo que permite el presupuesto?).
  \item \textbf{Ambiguo}: ``el sistema tiene que ser rápido y fácil de usar''. Mezcla dos atributos sin escala.
  Hay que separarlos (tiempo de respuesta; tiempo para sacar un turno, tasa de error de los usuarios mayores) y
  especificarlos con escala, test y niveles, como en b).
  \item También es una solución el ``recordatorio por SMS'': el requerimiento real es bajar el ausentismo; el SMS
  es una forma (podría ser WhatsApp o una llamada automática).
\end{itemize}

\paragraph{d) y e) Lista de comprobación y tabla de riesgos.}
\begin{center}\small
\begin{tabularx}{\textwidth}{@{}c X l l c l@{}}
\toprule
\# & Riesgo & Categoría & Conductor & Prob. & Impacto\\
\midrule
1 & Plazo de 8 meses insuficiente para web, app, SMS, lista de espera e integración & Proceso & Planif.\ ambiciosa & 70\% & 2 crítico\\
2 & Falta de experiencia del equipo en apps móviles & Plantilla & Tecnología poco familiar & 70\% & 2 crítico\\
3 & Pérdida de una de las 3 personas por la licencia (y reemplazo no previsto) & Plantilla & Personal no adecuado & 60\% & 2 crítico\\
4 & Resistencia de los administrativos de los centros & Negocio & --- & 60\% & 2 crítico\\
5 & Conectividad irregular en los centros rurales & Entorno / tecnología & --- & 60\% & 3 marginal\\
6 & Baja adopción por pacientes mayores & Negocio & Expectativas poco realistas & 50\% & 1 catastrófico\\
7 & Presupuesto fijo que vence a fin de año insuficiente & Tamaño / negocio & Subestimación del presupuesto & 50\% & 2 crítico\\
8 & Dificultades o demoras para integrarse con el padrón provincial & Tecnología & --- & 40\% & 2 crítico\\
9 & Requerimientos cambiantes (agenda, reglas de cancelación, lista de espera) & Cliente & Req.\ pobremente definidos & 40\% & 3 marginal\\
10 & Filtración de datos de salud & Tecnología / negocio & --- & 20\% & 1 catastrófico\\
11 & Pérdida de apoyo político por las elecciones & Negocio & Falta de soporte político & 20\% & 1 catastrófico\\
12 & Demora del proveedor de SMS o costo mayor al previsto & Entorno & --- & 20\% & 3 marginal\\
\bottomrule
\end{tabularx}
\end{center}
\textbf{Justificaciones} (las que se esperan en el examen): (1) el alcance es amplio y el plazo lo impone la
campaña de vacunación, que no se mueve: 70\%; crítico y no catastrófico porque puede salir una primera versión
reducida. (2) El enunciado dice que nadie hizo una app: es casi un hecho, pero se puede capacitar o hacer web
\emph{responsive}. (3) La licencia está prevista: lo incierto es si se consigue reemplazo; con 3 personas, perder
una es un tercio de la capacidad. (4) Los administrativos ya expresaron temor. (5) Ya existe el problema, pero solo
afecta a 3 de 12 centros y hay atención en mostrador: marginal. (6) Si los pacientes no lo usan, el problema no se
resuelve: catastrófico; 50\% porque hay mostrador y la población mayor es grande. (7) El presupuesto es fijo y
vence: si el proyecto se atrasa, se pierde. (8) Depende de otro organismo. (10) y (11) son graves pero poco
probables.

\textbf{Ordenamiento y línea de corte}: la tabla ya está ordenada por probabilidad y, a igual probabilidad, por
impacto (1, 2, 3, 4 tienen 70\% o 60\% y crítico; el 5 va después de los críticos del 60\% por ser marginal; el 6
va antes que el 7 por ser catastrófico). \textbf{Línea de corte en 50\%}: requieren mayor atención los riesgos 1 a
7. Casi todos nacen del mismo punto débil: \emph{mucho alcance, poco equipo y una fecha fija}.

\paragraph{f) RSGR de los tres primeros.}
\begin{itemize}
  \item \textbf{1. Plazo insuficiente}. \emph{Reducción}: priorizar con el Secretario una primera versión para
  la campaña (turnos web y mostrador, padrón y recordatorios) y dejar la app y la lista de espera para una
  segunda; planificación detallada con hitos y una reserva de tiempo. \emph{Supervisión}: reuniones de avance
  semanales; carta de tendencia de hitos. \emph{Gestión}: usar la reserva; mover funcionalidad a la segunda
  versión; pedir recursos adicionales.
  \item \textbf{2. Falta de experiencia en móvil}. \emph{Reducción}: capacitación temprana; empezar por una web
  \emph{responsive} y hacer una prueba de concepto de la app. \emph{Supervisión}: revisión técnica del prototipo en
  un hito fijo. \emph{Gestión}: contratar a un desarrollador móvil por tiempo acotado o dejar la app para la
  segunda versión.
  \item \textbf{3. Pérdida de personal por la licencia}. \emph{Reducción}: planificar la licencia en el
  cronograma; documentar en paquetes de trabajo lo que hace esa persona; capacitar a un compañero en sus tareas;
  pedir el reemplazo con anticipación. \emph{Supervisión}: seguimiento de la documentación y del pedido de
  reemplazo. \emph{Gestión}: reasignar tareas; incorporar un contratado; renegociar el alcance.
\end{itemize}

\subsection*{Ejercicio B2: cooperativa frutihortícola}
\paragraph{a) Requerimientos.} \textbf{Funcionales}: el productor publica semanalmente su mercadería (producto,
cantidad, precio, fecha de cosecha); el sistema da de baja los productos vencidos según los días de frescura de
su tipo; el cliente (comercio o particular) arma un pedido y elige día y franja de entrega; el sistema cobra con
tarjeta a través de la pasarela; el sistema emite la factura electrónica mediante el servicio del organismo
recaudador; el sistema genera la hoja de ruta diaria de cada repartidor; el repartidor confirma cada entrega
desde la app; el sistema calcula la liquidación mensual de cada productor (ventas menos comisión); el sistema
emite un reporte de ventas por producto y por productor. \textbf{No funcionales}: capacidad para triplicar los
pedidos en diciembre; usabilidad para productores sin experiencia digital; confiabilidad y exactitud de la
liquidación (hoy genera reclamos); seguridad de los pagos; disponibilidad de la app en ruta (posible mala
señal); interoperabilidad con la pasarela y con la facturación electrónica.

\paragraph{b) La planilla de Google.} Es una \textbf{solución disfrazada}. Preguntando \emph{¿para qué?} aparece
el requerimiento real: \emph{que los productores puedan informar su stock disponible de forma simple y que esa
información se use para vender, armar las entregas y liquidar}. La planilla es una solución posible (barata y
conocida), pero no da de baja los productos vencidos, no se integra con los pedidos ni con la liquidación y es
frágil ante errores de carga. Se trata como \emph{non-sancta}: quizás sirva para un primer piloto.

\paragraph{c) Riesgos (tabla ordenada).}
\begin{center}\small
\begin{tabularx}{\textwidth}{@{}c X l c l@{}}
\toprule
\# & Riesgo & Categoría & Prob. & Impacto\\
\midrule
1 & Productores que no cargan el stock o lo cargan mal (baja adopción) & Cliente / negocio & 70\% & 1 catastrófico\\
2 & Consultora sin conocimiento del dominio cooperativo (requerimientos mal entendidos) & Plantilla & 60\% & 2 crítico\\
3 & Reglas de liquidación y comisión ambiguas o cambiantes & Cliente & 60\% & 2 crítico\\
4 & Pérdida del apoyo al terminar el mandato de la presidenta & Negocio & 50\% & 1 catastrófico\\
5 & El sistema no soporta el pico de diciembre & Tecnología & 40\% & 2 crítico\\
6 & Demoras o problemas en la integración con la facturación electrónica & Tecnología & 40\% & 2 crítico\\
7 & Atrasos que impiden rendir los hitos del subsidio & Proceso & 40\% & 2 crítico\\
8 & Falta de señal en las rutas de reparto & Entorno & 40\% & 3 marginal\\
9 & Cambios en las comisiones de la pasarela de pagos & Negocio & 30\% & 3 marginal\\
10 & Fraude o filtración de datos de pago & Tecnología & 10\% & 1 catastrófico\\
\bottomrule
\end{tabularx}
\end{center}
\textbf{Línea de corte en 50\%}: riesgos 1 a 4. \textbf{RSGR}:
\textbf{1. Baja adopción de los productores}. \emph{Reducción}: involucrar a productores en el diseño; pantallas
muy simples; capacitación presencial; un referente de la cooperativa que cargue por quien no pueda.
\emph{Supervisión}: porcentaje de productores que cargan stock cada semana, contra lo esperado.
\emph{Gestión}: carga asistida por la administrativa; rediseño de la carga.
\textbf{2. Consultora sin conocimiento del dominio}. \emph{Reducción}: rondas de COS con la administrativa y
productores; prototipos validados; glosario. \emph{Supervisión}: validación de cada entrega con usuarios reales.
\emph{Gestión}: retrabajo priorizado; incorporar un referente de la cooperativa al equipo.

% =====================================================================
\subsection*{Ejercicio B3: funcional, no funcional o solución}
\begin{center}\small
\begin{tabularx}{\textwidth}{@{}c c X@{}}
\toprule
\# & Tipo & Comentario\\
\midrule
1 & F & función: está o no está.\\
2 & NF & atributo con escala (tiempo de respuesta) y condiciones de prueba.\\
3 & S & ¿por qué Oracle? Si es porque la empresa ya tiene licencias y DBA, el requerimiento real es una
\emph{restricción de diseño} o de costo; si no, es una solución que limita alternativas.\\
4 & F & salida con datos derivados (totales).\\
5 & S $+$ NF & el chatbot es una solución; el requerimiento real es reducir las llamadas (atributo: llamadas por
día, con nivel actual y planificado).\\
6 & NF & disponibilidad, con escala.\\
7 & mezcla & ``generar un manual'' es F; ``fácil de entender'' es NF (legibilidad): hay que separarlos.\\
8 & NF (restricción de diseño) & limita las soluciones posibles.\\
9 & F & función.\\
10 & NF & usabilidad (tiempo de aprendizaje) bien especificada.\\
\bottomrule
\end{tabularx}
\end{center}

% =====================================================================
\section{Bloque C --- Comparación de soluciones}
% =====================================================================
\subsection*{Ejercicio C1}
\textbf{Teoría}: resumen final §5.4.
\paragraph{a) Por impacto.} Sumas: S1 $= 60 + 15 + 10 + 35 = 120$; S2 $= 20 + 50 + 40 + 20 = \mathbf{130}$; S3
$= 30 + 30 + 35 + 25 = 120$. Se elige \textbf{S2}: mayor suma algebraica de impactos. Además es pareja (ninguno
por debajo de 20), mientras que S1 deja C en 10. Supuesto: los cuatro atributos pesan igual.

\paragraph{b) Atributo determinante.} Se excluye un atributo por vez:
\begin{center}\small
\begin{tabular}{lcccl}
\toprule
Sin & S1 & S2 & S3 & Preferida\\
\midrule
--- & 120 & 130 & 120 & S2\\
A & 60 & 110 & 90 & S2\\
B & \textbf{105} & 80 & 90 & \textbf{S1} (cambia)\\
C & \textbf{110} & 90 & 85 & \textbf{S1} (cambia)\\
D & 85 & 110 & 95 & S2\\
\bottomrule
\end{tabular}
\end{center}
\begin{resultado}{Resultado b}
Son determinantes \textbf{B y C}: la ventaja de S2 viene de cubrir 50\% de B y 40\% de C. Si se revisara la
importancia de cualquiera de los dos, ganaría S1.
\end{resultado}

\paragraph{c) Con costos.} Razón de impactos: S1 $= 120/30 = \mathbf{4{,}00}$; S2 $= 130/50 = 2{,}60$; S3 $=
120/40 = 3{,}00$. \textbf{Gana S1}: da más calidad por unidad de presupuesto.

\paragraph{d) Peor caso de C.} S1 aporta solo 10\% a C, por debajo del 15\% exigido: produce el \textbf{peor
caso} y se \textbf{elimina} (paso 3 del proceso). Entre S2 y S3, por razón de impactos gana \textbf{S3} (3,00
contra 2,60). Si no se consideraran costos, ganaría S2 (130).

\paragraph{e) Como líder de S2.} Legítimo: \textbf{reconsiderar la implementación} (la solución es
\emph{non-sancta}): bajar su costo (con 130 de impacto: por debajo del 32,5\% del presupuesto su razón supera el 4,00 de S1; en el escenario de d), con S1 descartada, alcanza con bajar del 43,3\% para superar el 3,00 de S3) o mejorar
A o D sin perjudicar el resto; o mostrar con evidencia (prototipo, medición) que su impacto real es mayor.
\textbf{Reconsiderar la relevancia de los atributos} solo es válido si se hace con el cliente y la gerencia, a
partir de los objetivos del POS, y no para acomodar el resultado: el líder de S2 no puede redefinir solo los
objetivos (los gerentes enuncian el problema; los \emph{problem solvers} buscan soluciones).
% =====================================================================
\section{Bloque D --- Estimación}
% =====================================================================
\textbf{Teoría}: resumen final §7.3.
\paragraph{D1 Delphi.} (a) Ordenadas: 9, 12, 13, 14, 15, 16, 22, 30. Con 8 valores, cada cuartil tiene 2: justifican
los del \textbf{cuartil inferior (9 y 12)} y los del \textbf{superior (22 y 30)}. La técnica lo pide porque en los
extremos suele estar la información que los demás no tienen (alguien ya hizo algo igual, o alguien ve un problema
técnico). (b) Promedio de la tercera pasada: $(14 + 15 + 15 + 16 + 17 + 15 + 16 + 18)/8 = 126/8 =
\mathbf{15{,}75}$ días. (c) Aprovecha el conocimiento del grupo y evita que mande el jefe o el más carismático.
Planning Poker es su versión ágil: también es grupal y discuten los extremos, pero estima en story points (relativos),
con cartas Fibonacci y revelación simultánea, y busca el consenso en vez de promediar.

\paragraph{D2 3 puntos.}
\begin{center}\small
\begin{tabular}{lcccc}
\toprule
Tarea & $E = (O + 4M + P)/6$ & & Rango $P - O$\\
\midrule
T1 & $(2 + 16 + 12)/6 = 5$ & & 10\\
T2 & $(5 + 32 + 17)/6 = 9$ & & 12\\
T3 & $(1 + 12 + 5)/6 = 3$ & & 4\\
T4 & $(10 + 48 + 26)/6 = 14$ & & 16\\
\bottomrule
\end{tabular}
\end{center}
(b) En T1 el pesimista (12) está mucho más lejos del más probable (4) que el optimista (2): la distribución es
asimétrica y la fórmula ``paga'' el riesgo de que salga mal. (c) La más incierta es \textbf{T4}: es la de mayor
distancia entre optimista y pesimista (16 días). T3 es simétrica: $E = M$.

\paragraph{D3 banda ancha.} Descartando mínimo y máximo de cada columna: $\bar O = (4 + 3 + 4)/3 = 3{,}67$; $\bar
M = (6 + 7 + 6)/3 = 6{,}33$; $\bar P = (10 + 12 + 11)/3 = 11$.
\begin{resultado}{Resultado D3}
$E = (3{,}67 + 4 \times 6{,}33 + 11)/6 = 40/6 \approx \mathbf{6{,}67}$ días. (Si en lugar de eso se descartara
solo el valor más extremo de cada columna, quedaría $\bar O = 4$, $\bar M = 6$, $\bar P = 10{,}5$ y $E \approx
6{,}42$: conviene aclarar en el examen qué criterio se usa.)
\end{resultado}

\paragraph{D4 duración y esfuerzo.} (a) $96 / (8 \times 1) = \mathbf{12}$ días. (b) Dedicación total $2 \times
0{,}6 = 1{,}2$: $96 / (8 \times 1{,}2) = \mathbf{10}$ días. (c) $96 / (8 \times 6) = \mathbf{2}$ personas full
time. (d) Por el \textbf{crash point}: cada persona agregada reduce cada vez menos la duración (coordinación,
comunicación, capacitación), y además muchas
tareas no se pueden partir. (e) El costo depende del \textbf{esfuerzo} (96 h-persona, con cualquier
distribución); la fecha de fin, de la \textbf{duración}.

% =====================================================================
\section{Bloque E --- Planificación}
% =====================================================================
\textbf{Teoría}: resumen final §7.5 y §7.6. Las tablas se calcularon con la convención de la cátedra y se
verificaron con un script que reproduce su resolución del ejercicio 3 del TP3. En los diagramas, arriba en azul
CTe y FTe; abajo en rojo CTa y FTa; en rojo las tareas y flechas críticas.

\subsection*{Ejercicio E1}
\textbf{Estrategia}: ordenar las tareas de modo que cada una venga después de sus predecesoras, ir hacia adelante
(máximo de los FTe $+1$), después hacia atrás (mínimo de los CTa $-1$), y calcular ventanas y márgenes.
\begin{center}
\begin{tikzpicture}[x=1.3cm, y=1.1cm]
  \tnode[crit]{A}{0,1}{A}{5}{1}{5}{1}{5}
  \tnode{B}{0,-1}{B}{3}{1}{3}{3}{5}
  \tnode{C}{2.5,1.6}{C}{4}{6}{9}{8}{11}
  \tnode[crit]{D}{2.5,0}{D}{6}{6}{11}{6}{11}
  \tnode{E}{2.5,-1.6}{E}{2}{4}{5}{14}{15}
  \tnode[crit]{F}{5,0.8}{F}{7}{12}{18}{12}{18}
  \tnode{G}{5,-1.6}{G}{3}{6}{8}{16}{18}
  \tnode[crit]{H}{7.5,0}{H}{4}{19}{22}{19}{22}
  \draw[dep] (A) -- (C.west); \draw[depc] (A) -- (D.west); \draw[dep] (B) -- (D.west); \draw[dep] (B) -- (E.west);
  \draw[dep] (C) -- (F.west); \draw[depc] (D) -- (F.west); \draw[dep] (E) -- (G);
  \draw[depc] (F) -- (H.west); \draw[dep] (G) -- (H.west);
\end{tikzpicture}
\end{center}
\begin{center}\small
\begin{tabular}{ccccccccl}
\toprule
Tarea & $d$ & CTe & FTe & CTa & FTa & $V$ & MT & ML\\
\midrule
A & 5 & 1 & 5 & 1 & 5 & 5 & 0 & 0\\
B & 3 & 1 & 3 & 3 & 5 & 5 & 2 & $\min(6, 4) - 3 - 1 = 0$\\
C & 4 & 6 & 9 & 8 & 11 & 6 & 2 & $12 - 9 - 1 = 2$\\
D & 6 & $\max(5, 3) + 1 = 6$ & 11 & 6 & 11 & 6 & 0 & 0\\
E & 2 & 4 & 5 & 14 & 15 & 12 & 10 & $6 - 5 - 1 = 0$\\
F & 7 & $\max(9, 11) + 1 = 12$ & 18 & 12 & 18 & 7 & 0 & 0\\
G & 3 & 6 & 8 & 16 & 18 & 13 & 10 & $19 - 8 - 1 = 10$\\
H & 4 & $\max(18, 8) + 1 = 19$ & 22 & 19 & 22 & 4 & 0 & $23 - 22 - 1 = 0$\\
\bottomrule
\end{tabular}
\end{center}
Hacia atrás, por ejemplo: FTa(B) $= \min(\text{CTa}_D, \text{CTa}_E) - 1 = \min(6, 14) - 1 = 5$, así que CTa(B) $= 3$.
\begin{resultado}{Resultado E1}
Duración: \textbf{22 días}. Camino crítico: \textbf{A -- D -- F -- H} ($5 + 6 + 7 + 4 = 22$). Los otros caminos:
A--C--F--H y B--D--F--H duran 20; B--E--G--H dura 12.
\end{resultado}
Detalle interesante: B tiene MT $= 2$ pero ML $= 0$, porque su sucesora E empieza justo al día siguiente; E tiene
MT $= 10$ y ML $= 0$ por lo mismo con G: el margen de 10 días lo comparten E y G.

\subsection*{Ejercicio E2}
Lo nuevo son las CC: la sucesora puede empezar el mismo día que la predecesora (CTe($p$) $\ge$ CTe($q$), sin $+1$).
\begin{center}\small
\begin{tabular}{ccccccccl}
\toprule
Tarea & $d$ & CTe & FTe & CTa & FTa & $V$ & MT & ML\\
\midrule
A & 2 & 1 & 2 & 3 & 4 & 4 & 2 & $5 - 2 - 1 = 2$\\
B & 4 & 1 & 4 & 1 & 4 & 4 & 0 & 0\\
C & 6 & $\max(2, 4) + 1 = 5$ & 10 & 5 & 10 & 6 & 0 & 0\\
D & 5 & CC con B: 1 & 5 & 12 & 16 & 16 & 11 & CC con E: $1 - 1 = 0$\\
E & 3 & CC con D: 1 & 3 & 12 & 14 & 14 & 11 & $4 - 3 - 1 = 0$\\
F & 8 & 11 & 18 & 11 & 18 & 8 & 0 & 0\\
G & 4 & 4 & 7 & 15 & 18 & 15 & 11 & $19 - 7 - 1 = 11$\\
H & 2 & 19 & 20 & 19 & 20 & 2 & 0 & 0\\
\bottomrule
\end{tabular}
\end{center}
Hacia atrás con CC: E tiene FTa $= \text{CTa}_G - 1 = 14$ y CTa $= 12$. D no tiene sucesoras FC, así que su FTa
inicial sería el fin del proyecto (20) y su CTa 16; pero su sucesora E (CC) exige CTa(D) $\le$ CTa(E) $= 12$, así
que CTa(D) $= 12$ y FTa(D) $= 16$.
\begin{resultado}{Resultado E2}
Duración: \textbf{20 días}. Camino crítico: \textbf{B -- C -- F -- H} ($4 + 6 + 8 + 2 = 20$). D y E tienen 11 días de
margen total pero 0 de margen libre: cada día que se atrasa D empuja a E, y E empuja a G.
\end{resultado}

\subsection*{Ejercicio E3 (los cuatro tipos de dependencias)}
\begin{center}
\begin{tikzpicture}[x=1.18cm, y=1.08cm]
  \tnode{A}{0,1.6}{A}{4}{1}{4}{6}{9}
  \tnode[crit]{B}{0,-1.2}{B}{6}{1}{6}{1}{6}
  \tnode{C}{2.4,1.6}{C}{5}{5}{9}{10}{14}
  \tnode[crit]{D}{2.4,-1.2}{D}{8}{7}{14}{7}{14}
  \tnode[crit]{E}{4.8,0.2}{E}{3}{15}{17}{15}{17}
  \tnode{F}{4.8,-2.8}{F}{4}{7}{10}{19}{22}
  \tnode{G}{7.2,1.8}{G}{7}{18}{24}{21}{27}
  \tnode[crit]{I}{7.2,-1.0}{I}{5}{18}{22}{18}{22}
  \tnode{H}{9.6,3.2}{H}{2}{1}{24}{30}{31}
  \tnode{L}{9.6,1.2}{L}{4}{25}{28}{28}{31}
  \tnode{J}{9.6,-0.9}{J}{6}{1}{18}{26}{31}
  \tnode[crit]{K}{9.6,-2.9}{K}{9}{23}{31}{23}{31}
  \draw[dep] (A) -- (C);
  \draw[depc] (B) -- (D);
  \draw[dep] (C) -- (E.west);
  \draw[depc] (D) -- (E.west);
  \draw[dep] (D.south west) |- node[pos=0.75, above, font=\tiny\bfseries, text=ambar]{CC} (F.west);
  \draw[dep] (E) -- (G.west);
  \draw[depc] (E) -- (I.west);
  \draw[dep] (G.east) to[out=0, in=0, looseness=1.6] node[pos=0.5, right, font=\tiny\bfseries, text=ambar]{FF} (H.east);
  \draw[dep] (G) -- (L.west);
  \draw[dep] (I.west) -- ++(-0.3,0) |- (6.9,0.05) -| node[pos=0.25, above, font=\tiny\bfseries, text=ambar]{CF} ([xshift=5mm]J.east) -- (J.east);
  \draw[depc] (I) -- (K.west);
  \draw[dep] (F) -- (K.west);
\end{tikzpicture}
\end{center}
Las flechas rotuladas salen del extremo de la predecesora que manda: la CC del comienzo de D, la FF del fin de G al
fin de H, la CF del comienzo de I al fin de J.

\paragraph{Planificación temprana.} A y B empiezan el día 1. C $=$ FTe(A) $+ 1 = 5$; D $=$ FTe(B) $+ 1 = 7$; E $=
\max(9, 14) + 1 = 15$ (manda D); F, CC con D: CTe $=$ CTe(D) $= 7$; G $=$ FTe(E) $+ 1 = 18$; I $= 18$; K $= \max(10,
22) + 1 = 23$; L $= 24 + 1 = 25$. \textbf{H (FF con G)}: nada restringe su comienzo, CTe $= 1$, FTe $= \max(1 + 2 -
1, \text{FTe}_G) = \max(2, 24) = 24$. \textbf{J (CF con I)}: CTe $= 1$, FTe $= \max(1 + 6 - 1, \text{CTe}_I) =
\max(6, 18) = 18$. El mayor FTe es el de K: \textbf{fin del proyecto día 31}.

\paragraph{Planificación tardía.} Sin sucesoras (H, J, K, L): FTa $= 31$. I: por K (FC) $23 - 1 = 22$; por J (CF)
CTa(I) $\le$ FTa(J) $= 31$, no restringe: FTa $= 22$, CTa $= 18$. G: por L $28 - 1 = 27$; por H (FF) FTa(G) $\le
31$, no restringe: FTa $= 27$, CTa $= 21$. F: por K $23 - 1 = 22$, CTa $= 19$. E: $\min(21, 18) - 1 = 17$, CTa
$= 15$. D: por E $15 - 1 = 14$; por F (CC) CTa(D) $\le$ CTa(F) $= 19$, no restringe: CTa $= 7$. C: $15 - 1 =
14$, CTa $= 10$. B: $7 - 1 = 6$, CTa $= 1$. A: $10 - 1 = 9$, CTa $= 6$.

\begin{center}\small
\begin{tabular}{cccccccccl}
\toprule
Tarea & $d$ & CTe & FTe & CTa & FTa & $V$ & MT $= V - d$ & ML & cálculo del ML\\
\midrule
A & 4 & 1 & 4 & 6 & 9 & 9 & 5 & 0 & $5 - 4 - 1$\\
B & 6 & 1 & 6 & 1 & 6 & 6 & 0 & 0 & $7 - 6 - 1$\\
C & 5 & 5 & 9 & 10 & 14 & 10 & 5 & 5 & $15 - 9 - 1$\\
D & 8 & 7 & 14 & 7 & 14 & 8 & 0 & 0 & $\min(15 - 14 - 1;\ \text{CC: } 7 - 7)$\\
E & 3 & 15 & 17 & 15 & 17 & 3 & 0 & 0 & $18 - 17 - 1$\\
F & 4 & 7 & 10 & 19 & 22 & 16 & 12 & 12 & $23 - 10 - 1$\\
G & 7 & 18 & 24 & 21 & 27 & 10 & 3 & 0 & $\min(\text{FF: } 24 - 24;\ 25 - 24 - 1)$\\
H & 2 & 1 & 24 & 30 & 31 & 31 & \textbf{29} & 7 & $32 - 24 - 1$\\
I & 5 & 18 & 22 & 18 & 22 & 5 & 0 & 0 & $\min(\text{CF: } 18 - 18;\ 23 - 22 - 1)$\\
J & 6 & 1 & 18 & 26 & 31 & 31 & \textbf{25} & 13 & $32 - 18 - 1$\\
K & 9 & 23 & 31 & 23 & 31 & 9 & 0 & 0 & $32 - 31 - 1$\\
L & 4 & 25 & 28 & 28 & 31 & 7 & 3 & 3 & $32 - 28 - 1$\\
\bottomrule
\end{tabular}
\end{center}
\begin{resultado}{Resultado E3}
Duración: \textbf{31 días}. Camino crítico: \textbf{B -- D -- E -- I -- K} ($6 + 8 + 3 + 5 + 9 = 31$), las tareas con
MT $= 0$ y temprana igual a tardía. Caminos que lo siguen: B--D--E--G--L $= 28$ (3 de margen) y A--C--E--I--K $=
26$.
\end{resultado}
\begin{aviso}{Igual que en la resolución de la cátedra}
En H y J el margen total \textbf{no} es FTa $-$ FTe ($31 - 24 = 7$ y $31 - 18 = 13$) sino $V - d$ ($29$ y $25$),
porque su CTe quedó en 1 y su FTe fue empujado por la FF y la CF. FTa $-$ FTe coincide, en cambio, con su margen
libre. Con la otra convención (la de MS Project: si se empuja el fin, se corre toda la tarea), H tendría CTe $=
23$ y MT $= 7$, y J CTe $= 13$ y MT $= 13$; el camino crítico y la duración no cambian.
\end{aviso}

\subsection*{Ejercicio E4: atrasos y compresión}
\begin{itemize}
  \item \textbf{D se atrasa 2 días}: D es crítica (MT $= 0$): el proyecto se atrasa 2 días y termina el día 24.
  \item \textbf{G se atrasa 8 días}: MT(G) $= 10$: el proyecto no se atrasa (G pasa a terminar el 16, antes de su FTa 18).
  \item \textbf{C se atrasa 3 días}: MT(C) $= 2$: se consumen 2 días de margen y el tercero atrasa el proyecto
  \textbf{1 día} (termina el 23). Ahora A--C--F--H dura 23 y pasa a ser el camino crítico.
  \item \textbf{E se atrasa 1 día}: ML(E) $= 0$, así que se atrasa el comienzo temprano de \textbf{G} (pasa del 6 al
  7); el proyecto no se atrasa porque E tiene MT $= 10$, que comparte con G.
  \item \textbf{Terminar en 20 días}: hay que acortar 2 días el camino crítico A--D--F--H. Pero A--C--F--H y
  B--D--F--H duran 20, así que conviene comprimir una tarea común a todos: \textbf{F} (o H). Opciones:
  agregar recursos a F (si no llegó a su crash point; cuesta dinero) o que F empiece con entregables parciales
  de D (FC $\to$ CC con lag; sube el riesgo de retrabajo). Acortando F en 2 días
  hay tres caminos críticos de 20 días: el proyecto queda más frágil.
\end{itemize}

% =====================================================================
\section{Bloque F --- Puntos de Función y COCOMO}
% =====================================================================
\textbf{Teoría}: resumen final §8.
\subsection*{Ejercicio F1}
\paragraph{a) PFNA.}
\begin{center}\small
\begin{tabular}{lccc}
\toprule
 & S4 & S5 & S6\\
\midrule
EE & $6 \times 3 = 18$ & $5 \times 4 + 2 \times 6 = 32$ & $8 \times 3 = 24$\\
SE & $5 \times 5 = 25$ & $4 \times 7 = 28$ & $3 \times 4 + 3 \times 7 = 33$\\
CE & $3 \times 3 + 2 \times 4 = 17$ & $6 \times 4 = 24$ & $5 \times 3 = 15$\\
AI & $4 \times 10 = 40$ & $2 \times 10 + 3 \times 15 = 65$ & $2 \times 7 = 14$\\
AE & $1 \times 5 = 5$ & $2 \times 7 = 14$ & $3 \times 10 = 30$\\
\midrule
\textbf{PFNA} & \textbf{105} & \textbf{163} & \textbf{116}\\
\bottomrule
\end{tabular}
\end{center}
\paragraph{b) FCT y PFA.} Se despliegan los rangos y se verifica que haya 14 factores:
\begin{center}\small
\begin{tabular}{l*{14}{c}c}
\toprule
 & F1 & F2 & F3 & F4 & F5 & F6 & F7 & F8 & F9 & F10 & F11 & F12 & F13 & F14 & TI\\
\midrule
S4 & 2 & 0 & 3 & 3 & 3 & 3 & 4 & 4 & 5 & 1 & 2 & 3 & 0 & 4 & 37\\
S5 & 3 & 3 & 4 & 4 & 4 & 4 & 4 & 5 & 5 & 3 & 1 & 2 & 1 & 5 & 48\\
S6 & 0 & 0 & 0 & 0 & 0 & 2 & 2 & 2 & 2 & 4 & 4 & 4 & 5 & 5 & 30\\
\bottomrule
\end{tabular}
\end{center}
Control por grupos: S4 $= 0 \cdot 2 + 1 \cdot 1 + 2 \cdot 2 + 3 \cdot 5 + 4 \cdot 3 + 5 \cdot 1 = 37$ (con $2 + 1 + 2 + 5 +
3 + 1 = 14$ factores); S5 $= 1 \cdot 2 + 2 \cdot 1 + 3 \cdot 3 + 4 \cdot 5 + 5 \cdot 3 = 48$; S6 $= 2 \cdot 4 + 4 \cdot 3 + 5
\cdot 2 = 30$.
\begin{resultado}{Resultado F1}
\begin{tabular}{lcccc}
 & PFNA & TI & FCT & PFA\\
S4 & 105 & 37 & $0{,}65 + 0{,}37 = 1{,}02$ & $105 \times 1{,}02 = \mathbf{107{,}10}$\\
S5 & 163 & 48 & 1,13 & $163 \times 1{,}13 = \mathbf{184{,}19}$\\
S6 & 116 & 30 & 0,95 & $116 \times 0{,}95 = \mathbf{110{,}20}$\\
\end{tabular}
\end{resultado}
(c) S6 tiene TI $= 30 < 35$: su entorno técnico es menos exigente que el promedio, así que el FCT es menor que 1 y
los PF bajan un 5\%. Con \textbf{TI $= 35$} (FCT $= 1$) los PFA son iguales a los PFNA. S6 tiene menos componentes
que S4, pero más PF, por sus 3 archivos externos complejos (30 puntos).

\subsection*{Ejercicio F2: contar desde la especificación}
\begin{center}\small
\begin{tabularx}{\textwidth}{@{}l c c c X@{}}
\toprule
Componente & Tipo & Compl. & Peso & Por qué\\
\midrule
Alta, modificación y baja de socios & 3 EE & simple & $3 \times 3 = 9$ & ingresan datos y mantienen el archivo de socios; cada operación cuenta\\
Registrar préstamo & EE & promedio & 4 & mantiene préstamos\\
Registrar devolución & EE & promedio & 4 & actualiza el préstamo\\
Comprobante de multa & SE & simple & 4 & tiene un dato \emph{derivado} (monto calculado)\\
Reporte mensual por categoría & SE & promedio & 5 & tiene totales\\
Consultar catálogo & CE & simple & 3 & muestra datos tal cual, sin modificar\\
Ficha del socio & CE & simple & 3 & ídem\\
Socios, Préstamos & 2 AI & simple & $2 \times 7 = 14$ & los mantiene el sistema\\
Libros y ejemplares & AI & promedio & 10 & ídem\\
Padrón de alumnos & AE & simple & 5 & lo mantiene el sistema académico; el nuestro solo lo lee\\
\bottomrule
\end{tabularx}
\end{center}
\begin{resultado}{Resultado F2}
PFNA $= 17 + 9 + 6 + 24 + 5 = \mathbf{61}$; con TI $= 40$, FCT $= 1{,}05$ y \textbf{PFA $= 64{,}05$}.
\end{resultado}
Casos dudosos: el comprobante de multa parece una consulta, pero calcula un monto: es salida. El padrón está en
la universidad, pero no lo mantiene nuestro sistema: es externo.

\subsection*{Ejercicio F3: COCOMO básico}
(a) $S = 120 \times 128 / 1000 = \mathbf{15{,}36}$ KLOC.
\begin{center}\small
\begin{tabular}{lcccc}
\toprule
Modo & $S^b$ & $E = a S^b$ (PM) & $D = 2{,}5 E^d$ (meses) & $P = E/D$\\
\midrule
Orgánico ($2{,}4$; $1{,}05$; $0{,}38$) & 17,61 & \textbf{42,3} & \textbf{10,4} & 4,1\\
Semi-embebido ($3{,}0$; $1{,}12$; $0{,}35$) & 21,32 & 64,0 & 10,7 & 6,0\\
Embebido ($3{,}6$; $1{,}20$; $0{,}32$) & 26,53 & \textbf{95,5} & \textbf{10,8} & 8,9\\
\bottomrule
\end{tabular}
\end{center}
(c) Del orgánico al embebido el esfuerzo se multiplica por 2,26 y la duración apenas un 4\%: el sistema embebido
no tarda mucho más, pero necesita más del doble de gente en paralelo ($d < 1$, y $d$ baja de 0,38 a 0,32). (d)
Liquidación de sueldos: \textbf{orgánico} (procesamiento de datos, requisitos negociables). Semáforos:
\textbf{embebido} (tiempo real, integrado con hardware, restricciones rígidas).

\subsection*{Ejercicio F4: COCOMO intermedio}
EAF $= 1{,}15 \times 1{,}30 \times 0{,}86 \times 1{,}07 \times 0{,}91 \times 1{,}23 = \mathbf{1{,}540}$. Esfuerzo
nominal intermedio: $3{,}2 \times 15{,}36^{1{,}05} = 56{,}35$ PM.
\begin{resultado}{Resultado F4}
$E = 56{,}35 \times 1{,}540 = \mathbf{86{,}8}$ personas-mes; $D = 2{,}5 \times 86{,}8^{0{,}38} = \mathbf{13{,}6}$
meses; unas 6,4 personas.
\end{resultado}
(b) Aumentan: RELY alta, CPLX muy alta (producto más exigente), LEXP baja (poca experiencia en el lenguaje) y SCED
muy baja. Reducen: ACAP alta (analistas capaces) y TOOL alta (buenas herramientas). SCED ``muy baja'' significa
un cronograma \emph{apretado} respecto del nominal: comprimir el plazo exige más gente y coordinación, y por eso
cuesta más esfuerzo (sin SCED, EAF $= 1{,}25$ y $E \approx 70{,}5$ PM).

\subsection*{Ejercicio F5: puntos objeto}
PO $= 8 \cdot 1 + 4 \cdot 2 + 2 \cdot 3 + 3 \cdot 2 + 2 \cdot 5 + 1 \cdot 8 + 3 \cdot 10 = 8 + 8 + 6 + 6 + 10 + 8 + 30
= \mathbf{76}$. NPO $= 76 \times 75/100 = \mathbf{57}$. Con PROD nominal (13): PM $= 57/13 \approx
\mathbf{4{,}4}$ personas-mes; con productividad baja (7): $57/7 \approx 8{,}1$.

% =====================================================================
\section{Bloque G --- Lanzamiento}
% =====================================================================
\textbf{Teoría}: resumen final §9.
\subsection*{Ejercicio G1: matriz persona/perfil}
\begin{center}\small
\begin{tabular}{l*{6}{c}r}
\toprule
Persona & Analista & DBA & Backend & Frontend & UX & Tester & \$/sem\\
\midrule
Ana Torres & \textcolor{critrojo}{\checkmark} & & & & & & 9.000\\
Bruno Díaz & & & \checkmark & & & & 11.000\\
Carla Gómez & & & \checkmark & & & & 10.500\\
Gabriela Sosa & & \textcolor{critrojo}{\checkmark} & \checkmark & & & & 12.000\\
Diego Paz & & & & \checkmark & & & 10.000\\
Hugo Vera & & & & \checkmark & \textcolor{critrojo}{\checkmark} & & 9.500\\
Elena Ruiz & & & & & & \checkmark & 7.500\\
Federico Lima & & & & & & \checkmark & 8.000\\
\midrule
Personas por perfil & \textbf{1} & \textbf{1} & 3 & 2 & \textbf{1} & 2 & \\
\bottomrule
\end{tabular}
\end{center}
\paragraph{b) Críticos.} \textbf{Ana} (única analista), \textbf{Gabriela} (única DBA) y \textbf{Hugo} (único UX): si
faltan, A1, A2 o A5 se detienen y no hay quién los reemplace. Matiz: como A4 y A5 se hacen en paralelo y Hugo tiene
que hacer A5, durante ese período \textbf{Diego} es el único frontend disponible: también es crítico en esas
semanas. Gabriela, Bruno y Carla no son críticos como backend (hay tres), pero Gabriela es la única DBA.

\paragraph{c) Asignación al menor costo.}
\begin{center}\small
\begin{tabular}{llll r}
\toprule
Actividad & Perfil & Candidatos & Asignado & \$/sem\\
\midrule
A1 relevamiento & analista & Ana & Ana & 9.000\\
A2 diseño de BD & DBA & Gabriela & Gabriela & 12.000\\
A3 API de turnos & backend & Bruno, Carla, Gabriela & Carla & 10.500\\
A7 notificaciones (paralela a A3) & backend & Bruno, Gabriela & Bruno & 11.000\\
A4 pantallas & frontend & Diego, Hugo & Diego (Hugo está en A5) & 10.000\\
A5 prototipos UX & UX & Hugo & Hugo & 9.500\\
A6 pruebas & tester & Elena, Federico & Elena & 7.500\\
\midrule
 & & & \textbf{Total} & \textbf{69.500}\\
\bottomrule
\end{tabular}
\end{center}
Para A3 y A7 hacen falta dos backend a la vez: se eligen los dos más baratos (Carla y Bruno) y Gabriela queda para
su perfil crítico. Para A4, Hugo es más barato que Diego, pero no puede hacer A4 y A5 al mismo tiempo.
\paragraph{d)} Paquetes de trabajo bien escritos de Ana, Gabriela y Hugo (el seguro ante su ausencia); capacitar a
otro (por ejemplo, Carla en bases de datos, Diego en UX); no asignarles tareas simultáneas; identificar de antemano
a quién contratar; tratarlo como un riesgo en el plan RSGR.

\subsection*{Ejercicio G2: asignación de horas (una solución válida)}
\begin{center}\small
\begin{tabular}{l*{5}{c}c}
\toprule
Tarea (semanas) & Laura & Martín & Sofía & Tomás & Rocío & Total\\
\midrule
T1 Gestión (1--8) & 64 & & & & & 64\\
T2 Relevamiento (1--2) & & 40 & & & 10 & 50\\
T3 Diseño (2--3) & & 20 & 30 & & & 50\\
T4 Inscripciones (3--6) & & 120 & & & & 120\\
T5 Reportes (3--6) & & & 100 & & & 100\\
T6 Pruebas (6--7) & & & & 60 & 10 & 70\\
T7 Capacitación (8) & & 30 & & & 8 & 38\\
\midrule
Total & 64 & 210 & 130 & 60 & 28 & 492\\
\bottomrule
\end{tabular}
\end{center}
\textbf{Carga semanal} (horas repartidas en forma pareja en la ventana de cada tarea): Laura 8 h todas las semanas;
Martín 20, 30, 40, 30, 30, 30, 0, 30; Sofía 0, 15, 40, 25, 25, 25, 0, 0; Tomás 30 en las semanas 6 y 7; Rocío 5,
5, 0, 0, 0, 5, 5, 8. \textbf{Nadie supera 40 h} (ni Rocío sus 10 h). La trampa: si Martín hacía solo el diseño, en
la semana 2 sumaba $20 + 25 = 45$ h. Se resolvió repartiendo T3 con Sofía (nivelación).
\textbf{Martín} (aspirante a LP) hace el relevamiento, parte del diseño y la capacitación: tareas de contacto con
el cliente y de visión global, que lo preparan para liderar (crecimiento individual); además es el respaldo de
Laura. Laura no programa: con 4 personas a cargo, gestiona y queda libre para imprevistos.

\subsection*{Ejercicio G3: pedido de cambio}
(a) El cliente lo pide \textbf{por escrito}; se abre una \textbf{ronda de COS} para entenderlo (¿qué medios de pago,
qué pasarela?); se revisa el pedido (se acepta para análisis); el equipo hace el \textbf{análisis de impacto} con
alternativas; el cliente elige; se actualiza el plan. (b) Alternativas: \emph{extender el cronograma 3 semanas}
(resultado 2: extensión); \emph{sumar un desarrollador contratado sin mover la fecha} (resultado 3: recursos
adicionales); \emph{mantener plazo y equipo y sacar otra funcionalidad de menor prioridad, que pasa a una segunda
versión} (resultado 5: revisión completa de los entregables). (c) \textbf{El cliente}: controla el alcance; el LP
le muestra el precio de cada alternativa.

\subsection*{Ejercicio G4: organización}
(a) La estructura I no sirve porque un LP no puede supervisar a más de 5 programadores. Con 12, la
\textbf{estructura II}: el LP coordina a 2 o 3 líderes de equipo, cada uno con 4 a 6 programadores (o la III, con
programadores senior a cargo de subgrupos). (b) El problema de la matricial: la persona tiene \textbf{dos jefes};
el gerente de línea es el permanente y su opinión suele pesar más que la del gerente de proyecto, que es temporario.
Funciona bien solo si ambos tienen igual autoridad. (c) Es \textbf{equipo contratado}: el tiempo del contrato no
varía (si la tarea se atrasa, se pierde el recurso), tiene que aprender el proyecto, su compromiso y la calidad
pueden ser un problema y requiere más supervisión.

% =====================================================================
\section{Bloque H --- Monitoreo y calidad}
% =====================================================================
\subsection*{Ejercicio H1: valor ganado}
(a) EV $= 20.000 + 15.000 + 25.000 = \mathbf{60.000}$ (regla 0/100: solo lo terminado).
(b) $SV = 60 - 70 = \mathbf{-10}$ mil; $SPI = 60/70 = \mathbf{0{,}857}$: va \textbf{atrasado} (hizo el 86\% de lo
planeado). $CV = 60 - 68 = \mathbf{-8}$ mil; $CPI = 60/68 = \mathbf{0{,}882}$: \textbf{sobrecosto} (cada peso
gastado produce 88 centavos de trabajo). Mirando solo PV y AC (70 contra 68) parecía ir bien.
(c) Contando T4 al 50\%, EV $= 70$: SPI $= 1$ y CPI $= 1{,}03$; el proyecto parecería al día y por debajo del costo.
Es la puerta del \emph{hope creep}: el ``50\%'' lo informa quien hace la tarea y suele ser optimista; contar solo
lo terminado es objetivo y verificable.

\subsection*{Ejercicio H2: carta de tendencia de hitos}
\textbf{P1}: después del 0, los hitos salen cada vez más tarde ($-2, -3, -5, -6$): \textbf{cuatro retrasos
sucesivos}: el proyecto se está yendo de control; hace falta una acción correctiva. \textbf{P2}: siete hitos
seguidos del mismo lado de la línea: \textbf{corridas sucesivas}: un cambio permanente en el cronograma
(tendencia), hay que replanificar. \textbf{P3}: un salto de $-1$ a $+8$: \textbf{cambio radical}: hay que
\textbf{renegociar el cronograma}.

\subsection*{Ejercicio H3: densidad y spoilage}
(a) A $= 30/12 = 2{,}5$; B $= 22/4 = \mathbf{5{,}5}$; C $= 18/20 = 0{,}9$ defectos/KLOC; sistema $= 70/36 \approx 1{,}94$.
Se revisa \textbf{B} (más del doble que A, aunque tenga menos defectos en número absoluto). (b) Spoilage $= 480/3.200 =
\mathbf{0{,}15}$: el 15\% del esfuerzo de desarrollo se fue en corregir fallas. (c) Defectos \emph{conocidos} no
son defectos \emph{existentes}: C puede tener pocos defectos porque se probó o se usó poco; además depende de cómo se
cuentan las LOC (si C tiene mucho código generado o verboso, su densidad baja artificialmente). Hay que comparar con
las mismas reglas y la misma ventana de uso.

\subsection*{Ejercicio H4: grafo del equipo}
$n = 6$; aristas: LP--A, LP--B, LP--C, LP--D, A--B, B--C, C--D, D--E: $e = 8$. Densidad $= 8/(6 \cdot 5/2) = 8/15
\approx 0{,}53$. Niveles individuales (grados): LP 4, A 2, B 3, C 3, D 3, E 1 (suma 16 $= 2e$). Promedio $= 16/6
\approx 2{,}67$. \textbf{Interpretación}: el LP concentra la comunicación (cuello de botella); E depende solo de
D: si D falta, E queda aislado. Agregar un canal E--LP o E--C reduciría ese riesgo.

\subsection*{Ejercicio H5: índice Fog}
Oraciones: 3. Palabras: $7 + 8 + 9 = 24$. Palabras de 3 o más sílabas: \emph{e-qui-po, en-tre-ga, sis-te-ma} (1.ª
oración); \emph{u-sua-rios, pan-ta-lla, du-ran-te, se-ma-na} (2.ª); \emph{pro-yec-to, a-prue-ba} (3.ª): 9 (no
cuentan \emph{marzo, prueban, cada, después, líder, versión, final}, de 2 sílabas).
\[ \text{Fog} = 0{,}4 \times \left(\frac{24}{3} + 100 \times \frac{9}{24}\right) = 0{,}4 \times (8 + 37{,}5) = \mathbf{18{,}2} \]
En la escala inglesa sería nivel de posgrado, pero el texto es simple: en castellano las palabras tienen más sílabas
y el Fog \textbf{infla} el resultado. Sirve para comparar versiones de un mismo texto, no como valor absoluto.

\subsection*{Ejercicio H6: riesgo de auditoría}
(a) $RD = RA/(RI \times RC) = 0{,}04/(0{,}8 \times 0{,}25) = \mathbf{0{,}20}$. (b) Con $RC = 0{,}5$: $RD =
0{,}04/0{,}4 = \mathbf{0{,}10}$. El auditor no controla RI ni RC; si los controles son peores, tiene que bajar su
riesgo de detección haciendo \textbf{más pruebas sustantivas} (muestras más grandes, más detalle).

\subsection*{Ejercicio H7: clasificaciones}
(a) (1) correctivo; (2) adaptativo; (3) perfectivo; (4) preventivo. (b) (1) preventivo (separación de funciones);
(2) detectivo (totales de control); (3) correctivo; (4) detectivo. (c) (1) abuso de privilegios; (2) acceso físico
ilegal; (3) virus y código malicioso; (4) hacking. (d) (1) portabilidad (transición); (2) mantenibilidad o
flexibilidad (revisión); (3) integridad (operación).

% =====================================================================
\section{Bloque I --- Dónde se responde cada pregunta conceptual}
% =====================================================================
\begin{center}\small
\begin{tabularx}{\textwidth}{@{}c X l@{}}
\toprule
\# & Tema & Resumen final\\
\midrule
1--2 & Proyecto, parámetros, etapas, clasificación & §3\\
3--4 & POS, meta/objetivo/criterio, FCE, COS & §4\\
5 & Riesgos: clasificaciones, checklist, conductores, RSGR, herramientas & §6\\
6 & Funcional, no funcional, solución; atributos & §5\\
7--12 & Actividad y tarea, crash point, estimación, WBS, diagramas, dependencias, planificaciones & §7\\
13--14 & Especificación, PF, COCOMO, COCOMO II & §8\\
15--16 & Organización del equipo, kick-off, cambios de alcance & §9\\
17--19 & Atributos externos, densidad, Fenton, McCall, defectos, mantenimiento, productividad & §11.1--11.2\\
20 & ISO & §11.3\\
21 & CMM y CMMI & §11.4\\
22 & Auditoría & §11.5\\
\bottomrule
\end{tabularx}
\end{center}
Para practicar estas preguntas en formato de opción múltiple, usá el quiz (\texttt{examen/quiz/index.html}).

\end{document}
